% YY11000C
% Immer durchzuführende Standardmaßnahmen
% 14.0
% 2013-11-30
:ref:`m-standardmassnahmen-immer`\label{topic:standardmassnahmen-immer}%
Die folgenden Maßnahmen sind grundsätzlich immer in einer *der Situation
angemessenen* Art und Weise durchzuführen. Im begründeten Ausnahmefall kann es
allerdings notwendig oder sinnvoll sein, dass manche Maßnahmen unterbleiben oder
aufgeschoben werden (Auf Grund des Selbstschutzes, \enquote{Aufklärung} eines bewusstlosen Patienten,
\ldots{}) oder angepasst werden müssen.

{\scriptsize (Die Reihenfolge ist der jeweiligen Situation anzupassen!)}
\begin{itemize}
  -   Einschätzungsblock (\prettyref{m:einschaetzungsblock}?
  -   Beengende oder behindernde Kleidung entfernen bzw.
  öffnen
  -   Situationsgerechte Lagerung
  -   *Wärmeerhalt* oder *Kühlung*
  -   Angemessene Erhebung der Anamnese 
  -   Angemessene Untersuchung 
  -   Spezielle Maßnahmen gemäß Verdachtsdiagnose(n)
  -   Patientenidentifkation
  -   Dokumentation, Aufklärung
  -   Psychischer Beistand
  -   Verlaufskontrolle, Patientenbeobachtung, Monitoring
%   \columnbreak
  -   Weiteres Vorgehen, je nach Bedarf und Situation:
  
  
 
  \begin{itemize}\label{tab:pam-weiterfuehrende}
    -   Transportentscheidung und ggfs. Transport an eine
    geeignete Einrichtung (Krankenanstalt).
    -   Notarzt-Nachforderung (bei Bedarf, z. B. Schmerztherapie,
    Aufklärung zwecks *Belassung*\index{Belassung} auf
    Patientenwunsch *trotz Behandlungsnotwendigkeit*, \ldots)
%     -   Unterlassung der Intervention wenn
%     *offensichtlich* kein sanitätsdienstliches Problem
%     besteht (Fehlalarmierung) 
%     Behandlungsbedarf/-wunsch
%     besteht und eine gesundheitliche Gefährdung des Patienten ausgeschlossen werden kann.
    \Ca{Rechtliche Hinweise beachten!}
    -   Patient begibt sich selbstständig in weitere
    Behandlung. Z. B. Bagatellverletzungen (z. B.
    oberflächliche Schürfwunde) ist es zulässig, dass sich
    der Patient selbständig in ärztliche Behandlung begibt.
    Das Anraten, eine ärztliche Behandlung aufzusuchen, ist
    zu dokumentieren und ggfs. vom Patienten per
    Unterschrift bestätigen zu lassen (Revers). Über Risiken
    muss aufgeklärt (und diese Aufklärung ebenfalls
    dokumentiert) werden.
  \end{itemize}
  -   Ggfs. Übergabe an weiterbehandelndes Personal
\end{itemize}

\Commentary[Immer durchzuführende Maßnahmen, Belassung]
  {Die Frage, wie nichtärztliches Personal
  bei einer Transportverweigerung oder vermeintlich nicht
  vorliegender Behandlungsbedürftigkeit/Spitalspflichtigkeit
  vorzugehen hat, ist nach wie vor nicht
  zufriedenstellend beantwortet.\ParBugfix
  Zur rechtlichen Lage in Österreich: Die Untersuchung auf
  das Vorliegen oder Nichtvorliegen von
  körperlichen und psychischen Krankheiten oder Störungen
  ist im \protect\S 2 Abs.~2 Z.~1 ÄrzteG geregelt und hat keine
  Entsprechung im SanG; somit fällt die Feststellung, ob
  eine Behandlungsbedürftigkeit besteht unter den
  Ärztevorbehalt, d. h. ist den Ärzten vorbehalten. \ParBugfix
  \protect\enquote{Ob Behandlungsbedarf vorliegt, kann nur
  von einem Arzt entschieden werden. Denn die dafür erforderliche
 Diagnose von (behaupteten) Krankheitszuständen fällt unter
 \protect\S 2 Abs.~2 Z.~1 ÄrzteG und ist daher den Ärzten
 vorbehalten.}} OGH (4 Ob 36/10p)
 \ParBugfix
 Das \enquote{selbstständige Begeben in ärztliche Behandlung}
 stellt keine Verneinung des Vorliegens einer Erkankung
 oder gesundheitlichen Störung sowie deren 
 grundsätzlicher Behandlungsnotwendigkeit dar.} % Commentary